\documentclass[11pt]{article}

\usepackage{../dad}

\begin{document}

\courseheader
  {Useful Data}
  {LMTH 2080 --- Data and Design with Python}
  {Day 01 \textbullet{} August 31, 2026}
  {Eugene Lang College of Liberal Arts}

\begin{keyidea}
Data is not useful simply because it exists. It becomes useful in relation to
a question, a problem, a person, and a decision.
\end{keyidea}

\section{What Would You Like to Know?}

Before we write any Python, begin with a problem or question.

Think about something you would genuinely like to understand, explain,
predict, compare, improve, or change. It can be large or small, serious or
ordinary.

\vspace{4pt}

\begin{dataprompt}
What is a problem or question that interests you?

\vspace{1.0in}
\end{dataprompt}

Now work backward from the question.

\begin{center}
\Large\bfseries
problem $\longrightarrow$ question $\longrightarrow$ evidence
$\longrightarrow$ data
\end{center}

\questionlabel

What would you need to know in order to investigate your question?

\vspace{0.7in}

\datalabel

What could you observe, measure, count, record, or collect that would provide
evidence?

\vspace{0.7in}

\section{A Data Hunt}

In a small group, choose \textbf{one} question to investigate. You have about
20 minutes to find real data on the web that might help.

You are not looking for an article that answers the question. You are looking
for the \emph{data} that would allow you to investigate it yourself.

Possible places to look include government open-data portals, research
organizations, public APIs, downloadable CSV or spreadsheet files, archives,
and datasets published by organizations.

\begin{checkpoint}
When you find a promising source, stop and inspect it before continuing to
search. The goal is not to collect links. The goal is to understand what you
have found.
\end{checkpoint}

For your best candidate, answer the following.

\begin{enumerate}[leftmargin=*,itemsep=10pt]
  \item \textbf{What are you trying to understand?}

  \vspace{0.45in}

  \item \textbf{Where did the data come from?} Who collected or published it?

  \vspace{0.45in}

  \item \textbf{What does one observation represent?} A person? A trip? A
  neighborhood? A day? A transaction? Something else?

  \vspace{0.45in}

  \item \textbf{What variables are available?} What has actually been
  measured or recorded?

  \vspace{0.45in}

  \item \textbf{What could you do with this data?} Name at least one question
  you could plausibly answer.

  \vspace{0.45in}

  \item \textbf{What is missing?} What would you like to know that this data
  cannot tell you?

  \vspace{0.45in}
\end{enumerate}

\begin{designprompt}
Who might find this data useful? What could they understand, decide, or do
with it?
\end{designprompt}

\newpage

\section{Data Is Made}

A dataset is not a neutral copy of the world. Someone decided what to count,
how to measure it, which categories to use, when to collect it, what to leave
out, and how to store it.

Return to the dataset your group found.

\limitationlabel

Who or what might be absent from this dataset?

\vspace{0.55in}

What decisions were made before you ever encountered the file?

\vspace{0.55in}

If you used this data to make a claim or a decision, what would you want
someone else to know about its limitations?

\vspace{0.55in}

\section{From Data on the Web to Data on Your Computer}

Eventually we want to do something with the data we find. That means knowing
where our files are and how our projects are organized.

Today we will use the \textbf{shell}: a text-based way of communicating with
your computer.

Open a terminal. Your instructor will provide a small course folder to work
with.

\begin{terminal}
pwd
\end{terminal}

\texttt{pwd} means \textbf{print working directory}. It answers a fundamental
question:

\begin{center}
\Large\bfseries Where am I?
\end{center}

Write the result here:

\vspace{0.4in}

Now ask the computer what is here.

\begin{terminal}
ls
\end{terminal}

What do you see?

\vspace{0.4in}

\subsection{Moving Around}

Directories are folders. Use \texttt{cd}---\textbf{change directory}---to
move into one.

\begin{terminal}
cd data
pwd
ls
\end{terminal}

To move back to the parent directory:

\begin{terminal}
cd ..
\end{terminal}

Try navigating around the course folder. At any moment, you should be able to
answer:

\begin{enumerate}[leftmargin=*]
  \item Where am I?
  \item What files and directories are here?
  \item How do I get where I want to go?
\end{enumerate}

\begin{keyidea}
You do not need to memorize shell commands today. You do need to begin
building a mental model of files, directories, and paths.
\end{keyidea}

\subsection{Making and Inspecting Things}

Try the following commands one at a time. After each command, use
\texttt{ls} and explain what changed.

\begin{terminal}
mkdir scratch
cd scratch
touch notes.txt
ls
cd ..
\end{terminal}

If the course folder contains a CSV file, we can inspect its first few lines
without opening a spreadsheet:

\begin{terminal}
head data/example.csv
\end{terminal}

A CSV is ultimately a text file. Soon Python will give us much more powerful
ways to work with it.

\newpage

\section{Keeping Track of Change with Git}

As we write code, analyze data, and revise projects, our files will change.
Git is a system for recording those changes.

Git answers another important question:

\begin{center}
\Large\bfseries What changed?
\end{center}

Move to the Git repository provided for today's exercise and ask Git about
its current state.

\begin{terminal}
git status
\end{terminal}

Make a small change to the file specified by your instructor. Then run:

\begin{terminal}
git status
\end{terminal}

What does Git notice?

\vspace{0.45in}

To tell Git that a change should be part of our next recorded version:

\begin{terminal}
git add <filename>
git status
\end{terminal}

To record the staged change:

\begin{terminal}
git commit -m "Add first notes"
\end{terminal}

And to see the history:

\begin{terminal}
git log --oneline
\end{terminal}

\begin{checkpoint}
A useful first mental model:

\begin{center}
files change $\longrightarrow$ Git notices $\longrightarrow$
we stage changes $\longrightarrow$ we commit a version
\end{center}
\end{checkpoint}

GitHub gives a Git repository a remote home. We will use Git and GitHub
throughout the semester to organize and share computational work. Branches,
merges, and other Git tools can wait.

\section{Before You Leave}

Return to the question you wrote at the beginning of class.

\questionlabel

What is one question you would genuinely like to investigate this semester?

\vspace{0.55in}

\datalabel

What data might help you investigate it?

\vspace{0.55in}

\limitationlabel

What might that data fail to capture?

\vspace{0.55in}

\section{For Next Time: Find Some Useful Data}

Before our next meeting, identify a problem or question that interests you and
find \textbf{two plausible data sources} that could help you investigate it.

For each source, record:

\begin{enumerate}[leftmargin=*]
  \item where the data comes from and who produced it;
  \item what one observation represents;
  \item several important variables;
  \item one question the data could help answer; and
  \item at least one important limitation or missing piece.
\end{enumerate}

Choose the more promising source and obtain a small sample of the actual data
if possible. Bring your question, your sources, and the data to our next
meeting.

You will also use the Git workflow from today to add a short Markdown file
describing what you found and make your first commit. Detailed submission
instructions will be provided with Assignment 0.

\begin{reflection}
You are not expected to leave today knowing Unix or Git. You should leave
with three habits:

\begin{enumerate}[leftmargin=*]
  \item begin with a question rather than a tool;
  \item ask whether the data you find is actually useful for that question;
  \item know where your files are and keep track of how your work changes.
\end{enumerate}
\end{reflection}

\end{document}
